\section{Introduction}
\label{sec:introduction}

\subsection{The Splitter We Need}

BISECT has several ways to split a graph node. \metis{} follows graph
boundary structure. \cvdgeo{} assigns tracts to geographic centers. \bfsg{}
grows two regions from seeds. \mka{} asks a different question:

\begin{quote}
Which straight sweep direction gives the best worst-half Reock score after the
node is divided into two population-balanced prefixes?
\end{quote}

The question is intentionally local. BISECT builds a final plan by recursive
bisection, so \mka{} is not a global optimizer over all $k$ districts. It is a
node-level structure primitive: take the current node, try many sweep angles,
choose the best two-way split under the implemented score, then continue down
the bisection tree.

\subsection{Show the Sweep, Not the Slogan}

The moving-knife idea is easy to say and easy to overstate. The useful visual
object is not merely a label such as ``47 degrees.'' It is the ordered prefix
created by projecting every tract onto the candidate axis.

\begin{center}
\begin{tabular}{c@{\qquad}c@{\qquad}c}
\textbf{Horizontal sweep} & \textbf{Diagonal sweep} & \textbf{Vertical sweep} \\
\begin{tabular}{|c|c|c|c|}
\hline
1 & 2 & 3 & 4 \\
\hline
1 & 2 & 3 & 4 \\
\hline
1 & 2 & 3 & 4 \\
\hline
\end{tabular}
&
\begin{tabular}{|c|c|c|c|}
\hline
1 & 1 & 2 & 3 \\
\hline
1 & 2 & 3 & 4 \\
\hline
2 & 3 & 4 & 4 \\
\hline
\end{tabular}
&
\begin{tabular}{|c|c|c|c|}
\hline
1 & 1 & 1 & 1 \\
\hline
2 & 2 & 2 & 2 \\
\hline
3 & 3 & 3 & 3 \\
\hline
\end{tabular}
\end{tabular}
\end{center}

The numbers are projection ranks. For each angle, the algorithm reads tracts
from rank 1 upward until the cumulative population reaches half of the node.
That prefix becomes one side of the candidate split; the remaining suffix
becomes the other side. The selected angle is the one whose two halves have the
largest minimum Reock score under the current centroid-MEC approximation.

\subsection{Why Reock Appears Here}

Reock compactness compares a district's area to the area of the smallest circle
that encloses it~\citep{reock1961}. It is a natural score for a sweep-based
splitter because bad directions create long, thin prefixes whose enclosing
circle grows quickly.

\begin{center}
\begin{tabular}{cc}
\textbf{Short enclosing circle} & \textbf{Long enclosing circle} \\
\begin{tabular}{|c|c|c|}
\hline
L & L & R \\
\hline
L & L & R \\
\hline
R & R & R \\
\hline
\end{tabular}
&
\begin{tabular}{|c|c|c|c|c|}
\hline
L & L & L & L & R \\
\hline
R & R & R & R & R \\
\hline
\end{tabular}
\end{tabular}
\end{center}

The left candidate is blocky. The right candidate stretches the left half
across most of the state, so its minimum enclosing circle is large relative to
its area. \mka{} prefers the former if population balance permits it.

\subsection{Implementation Status}

The implementation entry points live in the bisection runner:
\begin{center}
\small
\begin{tabular}{lll}
\texttt{split\_subgraph\_mka} &
\texttt{split\_subgraph\_mka\_direction} &
\texttt{run\_all\_splits\_mka}
\end{tabular}
\end{center}
The current contract is:

\begin{itemize}
\item centroid data are required; missing centroids produce a configuration
  error rather than a silent fallback;
\item coordinates are projected with the same Albers helper used by geographic
  CVD;
\item the default metric is Reock, computed from centroid points;
\item \texttt{polsby-popper} is parsed, but the current split routine falls
  back to Reock because perimeter data are not present there;
\item the post-sweep boundary rebalance is local and bounded, so final
  validation remains a required downstream step.
\end{itemize}

\subsection{Contributions}

This paper contributes an implementation-aligned account of \mka{}:

\begin{enumerate}
\item It explains the actual two-way recursive bisection routine and its
  orientation scan.
\item It gives visual checks for projection rank, population prefix, and
  worst-half Reock scoring.
\item It separates useful fair-division language from legal or empirical claims
  that require archived evidence packages.
\end{enumerate}

The result is a practical structure-layer paper, not a certification brief. It
shows what \mka{} currently does, why the choices are inspectable, and what
must be added before the method can support stronger publication claims.
